% !TeX root = ../main.tex
\chapter{Results}\label{ch:results}

\begin{guidance}
Report all the results obtained by applying the methods to the materials, usually in
tables and figures, in the same order as the experiments were described in
\autoref{ch:methods}. Describe what the results show objectively, without interpreting them
yet: the interpretation belongs to \autoref{ch:discussion} (some theses merge both chapters
into a single \enquote{Results and Discussion} chapter). Whenever possible, compare
different methods and include measures of dispersion and statistical significance.
\end{guidance}

\section{[First experiment]}\label{sec:experiment-1}

[Results of the first experiment. \autoref{tab:results} shows an example of a results table, in
which the best result of each column is highlighted in bold.]

\begin{table}[htbp]
  \centering
  \caption{Example of a results table: mean \textpm{} standard deviation over 10 runs. The
  best value of each column is shown in bold.}
  \label{tab:results}
  \begin{tabular}{@{}l c c r@{}}
    \toprule
    \textbf{Method} & \textbf{Accuracy} & \textbf{F1 score} & \textbf{Time (s)} \\
    \midrule
    Baseline        & 0.812 \textpm{} 0.014 & 0.774 \textpm{} 0.019 & 12.4 \\
    Method A        & 0.845 \textpm{} 0.011 & 0.801 \textpm{} 0.016 & 48.9 \\
    Method B        & 0.851 \textpm{} 0.009 & 0.815 \textpm{} 0.012 & 103.2 \\
    Proposed method & \textbf{0.873 \textpm{} 0.008} & \textbf{0.842 \textpm{} 0.010} & 61.7 \\
    \bottomrule
  \end{tabular}
\end{table}

\section{[Second experiment]}\label{sec:experiment-2}

[Results of the second experiment.]

\section{Comparison with the state of the art}\label{sec:comparison}

[Compare your results with those published in the literature reviewed in
\autoref{ch:background}.]
